% GENERATED FILE. Edit canonical lessons, then run npm run generate:books.
% canonical-source-hash: fnv1a64:538b39ca611bdadd
% canonical-lessons: LA-C77-mitto, LA-C77-maneo, LA-C77-cado, LA-C77-pono, LA-C77-invenio

\chapter{Send, Stay, Fall, Put, Find}
\label{ch:la-send-stay-fall-put-find}
\hlchaptermodality{77}

\begin{chapteropening}
\textbf{By the end of this chapter:} \emph{I can say to send, to stay, to fall, to put and to find.}

Complete the last lesson of chapter 77: I can say to send, to stay, to fall, to put and to find.
\end{chapteropening}

\section[mittō, mittere]{mittō, mittere — to send}

\label{lesson:LA-C77-mitto}



\begin{quote}
Before the new one: say the Latin for to swim, then the Latin for to call.
\end{quote}

\subsection*{You'll want to know: mittō, mittere}
\textbf{mittō, mittere} — \textquotedblleft{}to send\textquotedblright{}. English \emph{mission}.

\begin{cousinweb}
\textbf{mittō} has no secure origin; EtymDB 2.1 records a proposed Proto-Indo-European \emph{*meyth2-}, \textquotedblleft{}to exchange, to remove\textquotedblright{}, and Lewis \& Short call it doubtful. English \emph{admit} comes from it too.
\end{cousinweb}

\begin{grammarlens}[title={What you've built}]
One more word for this chapter.
\end{grammarlens}

\subsection*{Your turn}
\begin{itemize}
\raggedright
  \item \emph{Say it:} \emph{mittō}
  \item \emph{Say it:} \emph{mittō}, once more
  \item \emph{Say it:} say \emph{mittō}
  \item \emph{Recall:} say the Latin for to swim, then the Latin for to call, then say \emph{mittō} again
\end{itemize}

\begin{tcolorbox}[breakable,colback=teal!4,colframe=teal!35!black,title={Before you move on}]
What does mittō, mittere mean? (\textquotedblleft{}To send\textquotedblright{}.) Say it once more.
\end{tcolorbox}
\section[maneō, manēre]{maneō, manēre — to stay, to remain}

\label{lesson:LA-C77-maneo}



\begin{quote}
Before the new one: say the Latin for to call, then the Latin for to send.
\end{quote}

\subsection*{You'll want to know: maneō, manēre}
\textbf{maneō, manēre} — \textquotedblleft{}to stay, to remain\textquotedblright{}.

\begin{cousinweb}
\textbf{maneō} is inherited, usually traced to Proto-Indo-European \emph{*men-}, \textquotedblleft{}to stay\textquotedblright{} (EtymDB 2.1); Greek \emph{menō}, \textquotedblleft{}to stay\textquotedblright{}, is a cousin (Lewis \& Short). English \emph{remain} and \emph{mansion} come from it.
\end{cousinweb}

\begin{grammarlens}[title={What you've built}]
One more word for this chapter.
\end{grammarlens}

\subsection*{Your turn}
\begin{itemize}
\raggedright
  \item \emph{Say it:} \emph{maneō}
  \item \emph{Say it:} \emph{maneō}, once more
  \item \emph{Say it:} say \emph{maneō}
  \item \emph{Recall:} say the Latin for to call, then the Latin for to send, then say \emph{maneō} again
\end{itemize}

\begin{tcolorbox}[breakable,colback=teal!4,colframe=teal!35!black,title={Before you move on}]
What does maneō, manēre mean? (\textquotedblleft{}To stay, to remain\textquotedblright{}.) Say it once more.
\end{tcolorbox}
\section[cadō, cadere]{cadō, cadere — to fall}

\label{lesson:LA-C77-cado}



\begin{quote}
Before the new one: say the Latin for to send, then the Latin for to stay.
\end{quote}

\subsection*{You'll want to know: cadō, cadere}
\textbf{cadō, cadere} — \textquotedblleft{}to fall\textquotedblright{}.

\begin{cousinweb}
\textbf{cadō} is inherited, usually traced to Proto-Indo-European \emph{*kad-}, \textquotedblleft{}to fall\textquotedblright{} (EtymDB 2.1). English \emph{cadence}, \emph{case} and \emph{accident} come from it.
\end{cousinweb}

\begin{grammarlens}[title={What you've built}]
One more word for this chapter.
\end{grammarlens}

\subsection*{Your turn}
\begin{itemize}
\raggedright
  \item \emph{Say it:} \emph{cadō}
  \item \emph{Say it:} \emph{cadō}, once more
  \item \emph{Say it:} say \emph{cadō}
  \item \emph{Recall:} say the Latin for to send, then the Latin for to stay, then say \emph{cadō} again
\end{itemize}

\begin{tcolorbox}[breakable,colback=teal!4,colframe=teal!35!black,title={Before you move on}]
What does cadō, cadere mean? (\textquotedblleft{}To fall\textquotedblright{}.) Say it once more.
\end{tcolorbox}
\section[pōnō, pōnere]{pōnō, pōnere — to put, to place}

\label{lesson:LA-C77-pono}



\begin{quote}
Before the new one: say the Latin for to stay, then the Latin for to fall.
\end{quote}

\subsection*{You'll want to know: pōnō, pōnere}
\textbf{pōnō, pōnere} — \textquotedblleft{}to put, to place\textquotedblright{}. English \emph{position}.

\begin{cousinweb}
\textbf{pōnō} is contracted from an older \emph{*posinō} (Lewis \& Short), made of a prefix \emph{po-}, \textquotedblleft{}off, away\textquotedblright{}, and \emph{sinō}, \textquotedblleft{}to let, to leave\textquotedblright{}: to let something down in its place.
\end{cousinweb}

\begin{grammarlens}[title={What you've built}]
One more word for this chapter.
\end{grammarlens}

\subsection*{Your turn}
\begin{itemize}
\raggedright
  \item \emph{Say it:} \emph{pōnō}
  \item \emph{Say it:} \emph{pōnō}, once more
  \item \emph{Say it:} say \emph{pōnō}
  \item \emph{Recall:} say the Latin for to stay, then the Latin for to fall, then say \emph{pōnō} again
\end{itemize}

\begin{tcolorbox}[breakable,colback=teal!4,colframe=teal!35!black,title={Before you move on}]
What does pōnō, pōnere mean? (\textquotedblleft{}To put, to place\textquotedblright{}.) Say it once more.
\end{tcolorbox}
\section[inveniō, invenīre]{inveniō, invenīre — to find, to come upon}

\label{lesson:LA-C77-invenio}



\begin{quote}
Before the new one: say the Latin for to fall, then the Latin for to put.
\end{quote}

\subsection*{You'll want to know: inveniō, invenīre}
\textbf{inveniō, invenīre} — \textquotedblleft{}to find, to come upon\textquotedblright{}. English \emph{invent}.

\begin{cousinweb}
\textbf{inveniō} is built from \emph{in}, \textquotedblleft{}on, upon\textquotedblright{}, and \emph{veniō}, \textquotedblleft{}to come\textquotedblright{}: to come upon a thing (Lewis \& Short). English \emph{invent} first meant \textquotedblleft{}to find\textquotedblright{}.
\end{cousinweb}

\begin{grammarlens}[title={What you've built}]
One more word for this chapter.
\end{grammarlens}

\subsection*{Your turn}
\begin{itemize}
\raggedright
  \item \emph{Say it:} \emph{inveniō}
  \item \emph{Say it:} \emph{inveniō}, once more
  \item \emph{Say it:} say \emph{inveniō}
  \item \emph{Recall:} say the Latin for to fall, then the Latin for to put, then say \emph{inveniō} again
\end{itemize}

\begin{tcolorbox}[breakable,colback=teal!4,colframe=teal!35!black,title={Before you move on}]
What does inveniō, invenīre mean? (\textquotedblleft{}To find, to come upon\textquotedblright{}.) Say it once more.
\end{tcolorbox}
